% ===== 原文 pp.117–121（PDF p.127–131） =====
% [原文 p.117]
§9.2 所描述的玻色粘合条件必须用关于场 $\psi^I=(\psi^i,\hat\psi^i)$ 的费米子粘合条件来补充，即把费米子自由度总数的一半在缺陷线 $D$ 上置零；我们通过约束 $\psi^I$ 取入 $2n$ 维丛 $e^*TM$ 的一个 $n$ 维子丛——记作 $\cR_+$——来做到这一点。
% [Tutor Note: p.117 原文 “setting half of the total fermionic degrees of freedom to zero the defect line $D$” 缺了介词，应为 “to zero \emph{on} the defect line $D$”；译文按正确语法处理。另注意此处“取入 $n$ 维子丛 $\cR_+$”正是把一半费米子置零的实现方式：$e^*TM$ 是 $2n$ 维的，$\cR_+$ 是拉格朗日（$n$ 维）子丛]

确实，拓扑双膜 $(Y,F)$ 天然地配有一个中间维子丛，我们如下描述它：

\begin{definition}[7]
设 $(Y,F)$ 是命题 2、3、5 中所刻画的拓扑双膜。定义 $\cR_+$ 为 $Y$ 上如下丛：
\begin{equation}
\cR_+=\{u\in TY:\ G(u,w)=F(u,w)\ \ \forall w\in TY\}. \tag{9.2}
\end{equation}
\end{definition}

让我们看看这个子丛在前面讨论的两类特殊拓扑双膜中对应什么。当 $(Y,F)$ 是填满空间的情形时，$\cR_+$ 就是殆积结构 $G^{-1}F$ 的 $+1$ 特征丛。另一方面，对等距映射图像型双膜，这个子丛就是切丛 $TY$ 本身。

\begin{proposition}[6]
上面定义的子丛 $\cR_+$ 关于 $e^*TM$ 是拉格朗日的。特别地，它是 $n$ 维的。
\end{proposition}

给定满足几何拓扑双膜条件的 $(Y,F)$，我们施加前面写下的玻色粘合条件，并用“费米子取入 $\cR_+$”这一条件来扩充它们；此外，我们要求 $\cR_+$ 是 $Y$ 上的可积分布。为完整起见，我们在此记录完整的粘合条件组：
\begin{gather*}
\phi\in Y \\
G(v,w)=F(u,w) \\
G(s,w)=F(s,w) 
\end{gather*}
对所有 $w\in TY$ 成立，其中 $u^I\equiv\partial_0\phi^I$、$v^I\equiv\partial_1\phi^I$、$s^I\equiv\psi^I$。

现在我们来证明这些粘合条件定义了一个保持超对称的拓扑缺陷。我们写下带缺陷的 (0,1) 超对称 $\sigma$-模型的显式作用量（见 [57]）：
\begin{equation*}
\begin{aligned}
S=&\int_\Sigma \dd^2\sigma\left(-\tfrac12 g_{ij}\partial_\mu\phi^i\partial^\mu\phi^j-\tfrac12 b_{ij}\epsilon^{\mu\nu}\partial_\mu\phi^i\partial_\nu\phi^j+\tfrac{\ii}{2}g_{ij}\psi^i_+D_-\psi^j_+\right)\\
&+\int_{\hat\Sigma}\dd^2\sigma\left(-\tfrac12\hat g_{ij}\partial_\mu\hat\phi^i\partial^\mu\hat\phi^j-\tfrac12\hat b_{ij}\epsilon^{\mu\nu}\partial_\mu\hat\phi^i\partial_\nu\hat\phi^j+\tfrac{\ii}{2}\hat g_{ij}\hat\psi^i_+\hat D_-\hat\psi^j_+\right)+\int_D \dd\sigma^0\cA_I\partial_0\phi^I,
\end{aligned}
\end{equation*}
% [原文 p.118]
其中 $\partial_\pm=\partial_0\pm\partial_1$，而协变导数为
\begin{gather*}
D_\pm\psi^i_+=\partial_\pm\psi^i_++\partial_\pm\phi^j\Gamma^i_{jk}\psi^k_+ \\
\hat D_\pm\hat\psi^i_+=\partial_\pm\hat\psi^i_++\partial_\pm\hat\phi^j\hat\Gamma^i_{jk}\hat\psi^k_+ . 
\end{gather*}

变分此作用量并挑出定域在 $D$ 上的项，可得
\[
(\delta S)_D=\int_D \dd\sigma^0\left(-G_{IJ}\partial_1\phi^I\delta\phi^J+F_{IJ}\partial_0\phi^I\delta\phi^J-\frac{\ii}{2}G_{IJ}\psi^I_+\delta^{\text{cov}}\psi^J_+\right),
\]
其中 $G$ 与 $F$ 如前，并且我们引入了记号
\[
\delta^{\text{cov}}\psi^I_+=\delta\psi^I_++\delta\phi^J\Gamma^I_{JK}\psi^K_+
\]
表示协变变分；与 $\delta\psi^I$ 不同，它在靶空间坐标变换下协变地变换。这里 $\Gamma^I_{JK}$ 是 $X\times\hat X$ 上中性号差度规 $G$ 的 Levi-Civita 联络的系数。

$(\delta S)_D$ 中的前两项恒为零，因为玻色粘合条件与 §9.2 中的完全相同。此外，$(\delta S)_D$ 中的第三项可以借助下述命题证明为零：

\begin{proposition}[7]
设 $(M,G)$ 是（伪）黎曼流形，Levi-Civita 联络为 $\nabla$。设 $Y\subseteq M$ 是子流形，其世界体上配有闭 2-形式 $F$，并设 $\cR_+$ 是 (9.2) 中定义的子丛。若 $\cR_+$ 可积，则我们有
\[
(e^*\nabla)_u s\in\Gamma(\cR_+)
\]
对所有 $u\in\Gamma(TY)$ 与 $s\in\Gamma(\cR_+)$ 成立。这里 $e^*\nabla$ 是 $\nabla$ 在嵌入映射 $e:Y\to M$ 下的拉回。用文字说，子丛 $\cR_+$ 关于 $TY$ 是 $\nabla$-平行的。
\end{proposition}

\begin{proposition}[8]
项 $G_{IJ}\psi^I_+\delta^{\text{cov}}\psi^J_+$ 对所有允许变分为零。
\end{proposition}

注意在 $F=0$ 的情形，子丛 $\cR_+$ 等于 $TY$，上述命题就约化为如下令人惊讶的结果：$TY$ 的截面的协变导数仍在 $TY$ 中，即拉格朗日子流形自动是全测地的；这一结果也见于 [59]、[62]。

现在我们验证上述粘合条件定义了一个拓扑缺陷，并且保持理论的 (0,1) 超对称性。由 §9.2 的分析，能动张量的玻色部分自动光滑地粘合。剩下的要验证含费米子的贡献光滑地粘合；即
\[
(T^\mu{}_\nu)^f-(\hat T^\mu{}_\nu)^f=0
\]
% [原文 p.119]
对所有 $\mu,\nu$ 成立。这等价于条件
\[
G(\nabla_u s,s)=0,
\]
其中向量场 $u,s$ 如前定义。这个方程再次借助我们的平行性命题而得到满足。

最后，保持超对称性的条件是超流的 1-分量光滑地粘合。这等价于
\[
G(s,t)=0,
\]
其中 $s\equiv\psi_+$、$t\equiv\partial_+\phi$。这成立，依据是：由粘合条件，$\partial_+\phi$ 与 $\psi_+$ 都被约束为取入 $\cR_+$ 中的值，而 $\cR_+$ 是拉格朗日子丛。
% [Tutor Note: p.119 处“超流 1-分量粘合”指 $G(s,t)=0$：$s=\psi_+$ 与 $t=\partial_+\phi$ 都在 $\cR_+$ 中取值，而 $\cR_+$ 关于 $G$ 是拉格朗日的（即 $G|_{\cR_+}=0$，命题 6），故 $G(s,t)=0$ 自动成立。此处的“1-分量”与 p.109 的 $T^1{}_0$ 记号一致，均指 $\sigma^1$ 方向分量]

\section{(0,2) 超对称 \texorpdfstring{$\sigma$}{σ}-模型的拓扑缺陷}

现在考虑靶空间 $X$ 与 $\hat X$ 都是 Kähler 流形的情形，其 Kähler 形式为 $\omega$ 与 $\hat\omega$、复结构为 $j$ 与 $\hat j$。此时缺陷两侧的超对称 $\sigma$-模型都具有 (0,2) 超对称性；让我们从几何上分类那些保持两个超荷的拓扑双膜 $(Y,F)$。

我们施加上一节写下的粘合条件，并假定 (0,1) 超对称性的条件已满足。那么 (0,2) 超对称性将随之成立，只要沿 $D$ 有
\[
G'^1-\hat G'^1=0,
\]
其中 $G'$ 与 $\hat G'$ 是第二超流。这就是条件
\begin{equation}
\Omega(s,t)=0 \tag{9.3}
\end{equation}
% [Tutor Note: (9.3) 与命题 9 的条件 1 一致可核验：$\Omega(s,t)=G(Js,t)$；若 $J(\cR_+)=\cR_+$（条件 2）且 $\cR_+$ 关于 $G$ 拉格朗日（命题 6，即 $G|_{\cR_+}=0$），则对 $s,t\in\cR_+$ 有 $Js\in\cR_+$，故 $G(Js,t)=0$，即 (9.3) 成立——这正是“$R_+$ 是 $\Omega$-拉格朗日子丛”（条件 1）]
对所有 $s,t\in\cR_+$ 成立，其中
\[
\Omega=GJ
\]
是 $X\times\hat X$ 上关于复结构 $J=j\oplus\hat j$ 的 Kähler 形式。关于 (9.3) 何时成立，我们可以作如下陈述：

\begin{proposition}[9]
关于拓扑双膜 $(Y,F)$ 的下列条件是等价的：
% [原文 p.120]
\begin{enumerate}
\item $\cR_+$ 是 $\Omega$-拉格朗日子丛；
\item $J(\cR_+)=\cR_+$，即由子丛 $\cR_+$ 给出的叶状结构的叶是关于复结构 $J=j\oplus\hat j$ 的复子流形；
\item $Y$ 是 $X\times\hat X$ 的复子流形，且 $F$ 关于 $Y$ 上的诱导复结构具有类型 $(1,1)$；
\item 用广义几何的话说，$\cJ_J$ 作用在 $\tau^FY$ 的一个截面上给出 $\tau^FY$ 的另一个截面，其中
\[
\cJ_J=\begin{pmatrix}J&\\ &-J^T\end{pmatrix}
\]
是与复结构 $J$ 相应的广义复结构；
\item $(Y,F)$ 是关于复结构 $J$ 的 B-膜。
\end{enumerate}
\end{proposition}

% [Tutor Note: 命题 9 的五个条件互相等价可由 $\cR_G=-\cJ_\Omega\cJ_J$ 直接看出：两边左乘 $-\cJ_\Omega^{-1}=\cJ_\Omega$ 得 $\cJ_J=\cJ_\Omega\cR_G$，故任一被 $\tau^FY$ 保持的算子之积仍保持 $\tau^FY$——若 $\tau^FY$ 被 $\cR_G$ 与 $\cJ_\Omega$ 稳定化，则它也被 $\cJ_J$ 稳定化，反之亦然。这就是原文“any two ... then the third”的线性代数依据]
广义度规 $\cR_G$ 与广义复结构 $\cJ_\Omega$、$\cJ_J$ 之间由下式相联系：
\[
\cR_G=-\cJ_\Omega\cJ_J .
\]
这表明：若丛 $\tau^FY$ 被其中任意两个稳定化，则它自动被第三个稳定化。特别地，我们有

\begin{proposition}[10]
当双膜 $(Y,F)$ 关于辛形式 $\Omega=\omega\oplus-\hat\omega$ 是 A-膜、同时关于 $X\times\hat X$ 上的复结构 $J=j\oplus\hat j$ 是 B-膜时（在适当的可积性条件下），它是 (0,2) $\sigma$-模型的拓扑且保持超对称的缺陷。
\end{proposition}

\section{T-对偶}

我们简要评论我们的拓扑缺陷在实现靶空间为 $X$ 与 $\hat X$ 的玻色 $\sigma$-模型的对偶性中所起的作用。为具体起见，设 $X$ 与 $\hat X$ 是平坦环面，并设 $Y=X\times\hat X$ 是配有一个曲率为 $\cF=\dd\cA$ 的线丛的填满空间缺陷；我们按 §2 定义 2-形式 $F=-\cF-B$。

% [Tutor Note: p.120 原文此处作 “as in Section 2”，但按上下文应指 §9.2（2-形式 $F=-\cF-e^*B$ 在 §9.2、p.111 处定义）。原文的「Section 2」疑为笔误，译本照原文译作「§2」并在此说明。]

分隔理论 $X$ 与 $\hat X$ 的拓扑缺陷可以与分隔 $\hat X$ 与 $X'$ 的拓扑缺陷融合，得到分隔 $X$ 与 $X'$ 的拓扑缺陷。不可见缺陷（取对角线 $Y\subset X\times X$ 且 $F=0$）关于这一融合运算是幺元，而可逆缺陷是指能与另一个缺陷融合而得到不可见缺陷的缺陷。一般地，可逆的拓扑缺陷实现一个对偶性。\footnote{关于这一点在有理 CFT 语境中的近期讨论，见 [53]。}
% [Tutor Note: p.120 的“可逆缺陷实现对偶性”值得记下其确切含义：融合给出缺陷的幺半结构，不可见缺陷为幺元，可逆元即“逆缺陷”；把可逆缺陷插入配分函数相当于在两侧理论间作一次可逆变换，故它给出两个 $\sigma$-模型之间的等价。§9.6 随后把这一抽象陈述落实为环面上的 $O(n,n;\Z)$ 变换]

% [原文 p.121]
因此我们预期：若选取 2-形式 $\cF$ 使得 $(Y,F)$ 是可逆拓扑缺陷，则以 $X$ 为靶空间的 $\sigma$-模型将通过一个对偶变换与以 $\hat X$ 为靶空间的 $\sigma$-模型相联系。对环面，对偶变换群是 $O(n,n;\Z)$（即号差为 $(n,n)$ 的格的自同构群）[54]。

让我们固定背景数据 $(g,b)$ 并取
\[
\cF=\begin{pmatrix}f&h\\ -h^T&\hat f\end{pmatrix},
\]
其中 $f$、$\hat f$ 是反对称 $n\times n$ 矩阵。由于 $\cF$ 是环面 $X\times\hat X$ 上线丛的曲率，其元素被约束为整数。一般地，非对角矩阵 $h$ 的行列式可以取任意整数值；不过，让我们现在通过假定 $\det h=\pm1$（或等价地 $h\in GL(n,\Z)$）来限制 $\cF$ 的形式。

对填满空间缺陷，拓扑缺陷条件为 $(G^{-1}F)^2=1$；按分块写开，它等价于 $X$ 的背景数据 $(g,b)$ 与 $\hat X$ 的背景数据 $(\hat g,\hat b)$ 之间的如下关系：
\begin{gather*}
\hat g=h^T\left(g-(b+f)g^{-1}(b+f)\right)^{-1}h \\
0=(b+f)g^{-1}h+h\hat g^{-1}(\hat b-\hat f). 
\end{gather*}

我们认出这就是如下对偶变换：先把 B-场平移矩阵 $f$，接着作一次 T-对偶（即把数据 $g+b$ 取逆），再作一次生成环面的格之基变换（由幺模矩阵 $h$ 参数化），最后再作一次 B-场平移（这次平移 $\hat f$）（见 [54] 的 (2.4.25)、(2.4.26)、(2.4.39) 式）。相应的对偶元素为
\[
\begin{pmatrix}\hat f h^{-1}&h^T+\hat f h^{-1}f\\ h^{-1}&h^{-1}f\end{pmatrix}\in O(n,n;\Z).
\]
（取 $f=\hat f=0$、$h=1$ 的类似计算在 [51] 中作出。）从这一表达式可以明显看出，$\cF$ 元素的量子化蕴含该对偶元素各元素的量子化。
% [Tutor Note: p.131 的对偶元素可用“B-平移—T-对偶—基变换—B-平移”四步复合核对：T-对偶在 $O(n,n)$ 中由 $\begin{pmatrix}0&1\\1&0\end{pmatrix}$ 型矩阵实现，B-平移由 $\begin{pmatrix}1&0\\ \theta&1\end{pmatrix}$ 实现，基变换由 $\diag(h,h^{-T})$ 实现；按此顺序相乘即得上式，且幺模性保证它落在 $O(n,n;\Z)$ 中。另核对：$h\to h^{-1}$ 的出现与 T-对偶把 $g+b$ 取逆一致]
